\documentclass{article}
\usepackage{style}
\usepackage{graphicx} % Required for inserting images
\newtheorem{thm}{Theorem}
\newtheorem{prop}{Proposition}
\theoremstyle{remark}
\newtheorem*{rmk}{Remark}
\usepackage{hyperref}
\newenvironment{proofSketch}{\noindent \textbf{Proof (sketch):}}
{\hfill\\}

\title{Hill Function Learning}
\setlength{\parindent}{0pt}

\begin{document}

\maketitle
\section{Setup}
Let $\Omega_S$ be a space of model inputs (RNA sequences, primary and secondary structures, etc.).
\newcommand{\distr}{\rho}

There exists $f_{H}:\Omega_S\to\RR^4$ with $f_{H}(s)=(Y_{\min}(s),Y_{\max}(s),K_{1/2}(s),n(s))$ parameterizing the Hill curve
\begin{align*}
    H(x;Y_{\min},Y_{\max},K_{1/2},n)
     & = Y_{\min} + \frac{Y_{\max} - Y_{\min}}{1 + \left(\frac{K_{1/2}}{x}\right)^n} \\
     & = Y_{\min} + (Y_{\max} - Y_{\min})\sigma\left(n\log\frac{x}{K_{1/2}}\right)   \\
     & = H(x;f_{H}(s)) = H(x;s),
    \intertext{and the fold value}
    \mathrm{fold}(s) & = \frac{H(1;s)}{H(0;s)}.
\end{align*}
The predictor is $\mu_{\theta}:\Omega_S\to\RR^4$ with $\mu_{\theta}(s)\approx f_{H}(s)$.

\section{Data}
Each input $S$ comes with a finite set of noisy observations
\[
    Z = \{(x_j, y_j)\}_{j=1}^{m(S)} \subset \RR \times \RR,
    \qquad (S,Z)\sim\distr,
\]
where $(0,\cdot)$ and $(1,\cdot)$ are always present and interior concentrations may appear.
Each pair satisfies $y = H(x;f_H(S)) + \varepsilon$ with $\expect{\varepsilon\mid S,x}=0$ and independent $\varepsilon$ across points.

\section{Baseline-normalized loss}
Fold is a relative quantity, so normalize by the predicted basal level $H(0;\mu_{\theta}(S))$:
\begin{align*}
    \widetilde{H}(x;\mu_{\theta}(S))
     & \defeq \frac{H(x;\mu_{\theta}(S))}{H(0;\mu_{\theta}(S))},
    \qquad
    \widetilde{H}(0;\mu_{\theta}(S)) = 1,
    \qquad
    \widetilde{H}(1;\mu_{\theta}(S)) = \mathrm{fold}_\mu(S).
\end{align*}
The per-point kernel is
\begin{align*}
    \ell(y,x;\mu_{\theta}(S))
     & \defeq
    \left(
        \frac{y}{H(0;\mu_{\theta}(S))}
        - \widetilde{H}(x;\mu_{\theta}(S))
    \right)^2.
\end{align*}
Each input contributes the mean over its observation set,
\begin{align*}
    \ell_{\theta}(S,Z)
     & \defeq \frac{1}{|Z|}\sum_{(x,y)\in Z} \ell(y,x;\mu_{\theta}(S)),
\end{align*}
so $|Z|$ does not weight sequences with more concentrations more heavily.
The population loss is $\mathcal{L}(\theta)\defeq\expect[(S,Z)]{\ell_{\theta}(S,Z)}$, which under independence of $\varepsilon$ from the model error decomposes as
\begin{align*}
    \mathcal{L}(\theta)
     & = \expect[(S,Z)]{
        \frac{1}{|Z|}\sum_{(x,y)\in Z}
        \left(
            \frac{H(x;f_H(S))}{H(0;\mu_{\theta}(S))}
            - \widetilde{H}(x;\mu_{\theta}(S))
        \right)^2
    }
    + \expect[(S,Z)]{
        \frac{1}{|Z|}\sum_{(x,y)\in Z} \frac{\varepsilon^2}{H(0;\mu_{\theta}(S))^2}
    }.
\end{align*}
When $Z=\{(0,y_0),(1,y_1)\}$,
\begin{align*}
    \ell_{\theta}(S,Z)
     & = \tfrac12\left(\frac{y_0}{H(0;\mu_{\theta}(S))} - 1\right)^2
    + \tfrac12\left(\frac{y_1}{H(0;\mu_{\theta}(S))} - \mathrm{fold}_\mu(S)\right)^2.
\end{align*}
Interior points enter through the same $\ell$ with weight $1/|Z|$ and constrain $K_{1/2}$ and $n$ via $\widetilde{H}$.

\section{Constraints}
Enforce $Y_{\min}(S)>0$ and $H(x;\mu_{\theta}(S))>0$ (e.g.\ softplus outputs with a floor).
Require a minimum dynamic range $Y_{\max}\ge Y_{\min}+\delta$ for $\delta>0$ so the curve cannot collapse to flat.
Since $Y_{\min}=H(0;\mu_{\theta})$ is learned and appears in every denominator, basal and fold are estimated jointly from the same normalization.

\end{document}
